\documentclass[10pt,a4paper]{article}
\usepackage{amsmath,amssymb,amsthm,graphicx}
% Minimal retained preamble for lane-x1 cards taken from
% sources/sigma-2023-004-law-of-large-numbers-for-roots-of-finite-fre/pinned-original-source.tex
% (doi:10.3842/SIGMA.2023.004, SIGMA 19 (2023), article 004, CC BY 4.0).
%
% Retained verbatim from the source preamble: the theorem declarations
% under the per-section shared counter and the per-section equation
% numbering.
%
% Omitted and why: the \documentclass{sigma} lines (texlive carries no
% sigma.cls; per the section 8.8.6x precedent the assembler supplies the
% article baseline class) and the SIGMA front-matter commands, which belong
% to the journal layout and are not used by the retained spans.
\usepackage{amsmath,amssymb,amsthm,mathrsfs,mathtools,graphicx}

\numberwithin{equation}{section}

\newtheorem{Theorem}{Theorem}[section]
\newtheorem{Corollary}[Theorem]{Corollary}
\newtheorem{Lemma}[Theorem]{Lemma}
\newtheorem{Proposition}[Theorem]{Proposition}
 { \theoremstyle{definition}
\newtheorem{Definition}[Theorem]{Definition}
\newtheorem{Note}[Theorem]{Note}
\newtheorem{Example}[Theorem]{Example}
\newtheorem{Remark}[Theorem]{Remark}
\newtheorem{Conjecture}[Theorem]{Conjecture}}

\title{Law of large numbers for roots of finite free multiplicative convolution powers}
\author{Fujie, Katsunori, Ueda, Yuki}
\date{doi:10.3842/SIGMA.2023.004; CC BY 4.0}
\begin{document}
\maketitle
\setcounter{section}{3}
\setcounter{subsection}{0}
\setcounter{Theorem}{1}
\setcounter{equation}{1}
\noindent\textit{Editorial note.} All delivered mathematical text is the author's own, copied from the pinned source file sources/sigma-2023-004-law-of-large-numbers-for-roots-of-finite-fre/pinned-original-source.tex; no statement or proof step was rewritten, reordered or removed. The counted span is the author environment and its own proof environment. References inside the span are handed to the assembler as navigation stubs; whether each of them is load-bearing is recorded on the stub itself.\par\medskip
\noindent\textit{Editorial scope.} Single Theorem card cut from the SIGMA article doi:10.3842/SIGMA.2023.004 (sigma-2023-004-law-of-large-numbers-for-roots-of-finite-fre); statement and proof are the author own text at source lines 253-260 and 266-309; the context spans are the author's own notation and definition passages at 112-120, 195-204, 225-231.\par\medskip
\noindent\textit{Allowed standard prerequisites (editorial).}\begin{itemize}
\item author notation block (elementary symmetric functions and their normalized variants) (context span 112-120). carried as a context span. supplies the notation the statement and proof use
\item author Definition defn (finite free multiplicative convolution) (context span 195-204). carried as a context span. the convolution operation the lemma is about; kept as context so the label defn resolves in-card
\item author Lemma lem:free\_multiplicative\_polynomial statement (supplies eq:symm\_root) (context span 225-231). carried as a context span. the convolution-power identity eq:symm\_root the proof applies; kept as context so the labels resolve in-card
\end{itemize}\medskip
\begin{Theorem}\label{thm:LLN}
Consider a monic polynomial $p$ of degree $d$ with non-negative real roots $\Lambda$ and let $k=k(p)$ be the number of zeros in $\Lambda$.
Let $\Lambda^{(n)}:=\big\{\lambda_{1}^{(n)} \ge \lambda_{2}^{(n)} \ge\dots \ge \lambda_{d}^{(n)} \big\} $ be the set of non-negative real roots of $p^{\boxtimes_d n}$.
Then we get
\begin{align*}
\lim_{n\rightarrow\infty}\big(\lambda_{i}^{(n)}\big)^{\frac{1}{n}} =\frac{\widetilde{e}_i(\Lambda)}{\widetilde{e}_{i-1}(\Lambda)}, \qquad 1 \le i \le d-k.
\end{align*}
\end{Theorem}

For a finite multiset $\Lambda=\{\lambda_1,\dots, \lambda_d\}$ of complex numbers, the $i$-th elementary symmetric polynomials $e_i(\Lambda)$ is denoted by
\begin{align*}
e_i(\Lambda):= \sum_{J \subset [d],\, |J|=i} \bigg( \prod_{j \in J} \lambda_j \bigg), \qquad e_0(\Lambda):=1,
\end{align*}
where $[d]:=\{1,2,\dots, d\}$.
In addition, we define
\begin{align*}
\widetilde{e}_i(\Lambda):=\frac{e_i(\Lambda)}{\binom{d}{i}}.
\end{align*}

\begin{Definition}\label{defn}
For monic polynomials $p$ and $q$ of degree $d$ which have only non-negative real roots:
\begin{align*}
p(x)=\sum_{i=0}^d (-1)^i p_ix^{d-i}, \qquad \text{ and } \qquad q(x)=\sum_{i=0}^d (-1)^i q_ix^{d-i},
\end{align*}
the {\it finite free multiplicative convolution} $p\boxtimes_d q$ is defined by
\begin{align*}
\big(p\boxtimes_d q\big) (x):= \sum_{i=0}^d(-1)^i \frac{p_iq_i}{\binom{d}{i}} x^{d-i}.
\end{align*}
\end{Definition}

\begin{Lemma}\label{lem:free_multiplicative_polynomial}
Let $p$ be a monic polynomial of degree $d$ with non-negative real roots $\Lambda$. Then we have
\begin{align} \label{eq:symm_root}
\widetilde{e}_{i}\big(\Lambda^{(n)}\big)= \widetilde{e}_{i}(\Lambda)^n, \qquad 0\le i \le d.
\end{align}
In particular, the number of zeros in $\Lambda$ is the same as the one in $\Lambda^{(n)}$.
\end{Lemma}

\begin{proof}
For $1\le i \le d-k$, the equation \eqref{eq:symm_root} implies that
\begin{align}\label{eq:Lambda_S}
\frac{\widetilde{e}_{i}\big(\Lambda^{(n)}\big)}{\widetilde{e}_{i-1}\big(\Lambda^{(n)}\big)}= \left( \frac{\widetilde{e}_{i}(\Lambda)}{\widetilde{e}_{i-1}(\Lambda)}\right)^n.
\end{align}

For $i=1$, the equation \eqref{eq:Lambda_S} implies that
\begin{align*}
\frac{\lambda_{1}^{(n)}+\dots + \lambda_{d}^{(n)}}{d}=\widetilde{e}_1\big(\Lambda^{(n)}\big)=\widetilde{e}_1(\Lambda)^n.
\end{align*}
Since
\begin{align*}
\frac{\lambda_{1}^{(n)}}{d} \le \frac{\lambda_{1}^{(n)}+\dots + \lambda_{d}^{(n)}}{d} \le \lambda_{1}^{(n)},
\end{align*}
we obtain
\begin{align*}
\widetilde{e}_1(\Lambda) \le \big(\lambda_{1}^{(n)}\big)^{\frac{1}{n}} \le d^{\frac{1}{n}}\widetilde{e}_1(\Lambda),
\end{align*}
and therefore {$\big(\lambda_{1}^{(n)}\big)^{\frac{1}{n}} \rightarrow \widetilde{e}_1(\Lambda)$ as $n\rightarrow\infty$.}

For {$2\le i \le d-k$}, we have
\begin{align*}
\frac{e_{i}\big(\Lambda^{(n)}\big)}{e_{i-1}\big(\Lambda^{(n)}\big)} &= \frac{\sum_{J \subset [d],\ |J|=i} \left( \prod_{j \in J} \lambda_{j}^{(n)} \right)}{ \sum_{J \subset [d],\ |J|=i-1} \left( \prod_{j \in J} \lambda_{j}^{(n)} \right)}\ge \frac{\sum_{J \subset [d],\ |J|=i} \left( \prod_{j \in J} \lambda_{j}^{(n)} \right)}{\binom{d}{i-1} \lambda_{1}^{(n)}\lambda_{2}^{(n)}\cdots \lambda_{i-1}^{(n)}}\\
&\ge \frac{\lambda_{1}^{(n)}\lambda_{2}^{(n)}\cdots \lambda_{i}^{(n)}}{\binom{d}{i-1} \lambda_{1}^{(n)}\lambda_{2}^{(n)}\cdots \lambda_{i-1}^{(n)}}=\binom{d}{i-1}^{-1} \lambda_{i}^{(n)}.
\end{align*}
Similar to the estimation above, we obtain
\begin{align*}
\frac{e_{i}\big(\Lambda^{(n)}\big)}{e_{i-1}\big(\Lambda^{(n)}\big)}  &\le \frac{\binom{d}{i}\lambda_{1}^{(n)}\lambda_{2}^{(n)}\cdots \lambda_{i}^{(n)} }{\lambda_{1}^{(n)}\lambda_{2}^{(n)}\cdots \lambda_{i-1}^{(n)}} = \binom{d}{i} \lambda_{i}^{(n)}.
\end{align*}
Consequently, we have
\begin{align*}
\binom{d}{i}^{-1} \frac{e_i\big(\Lambda^{(n)}\big)}{e_{i-1}\big(\Lambda^{(n)}\big)} \le \lambda_{i}^{(n)} \le \binom{d}{i-1} \frac{e_i\big(\Lambda^{(n)}\big)}{e_{i-1}\big(\Lambda^{(n)}\big)},
\end{align*}
and therefore
\begin{align*}
\binom{d}{{i-1}}^{-1} \left(\frac{\widetilde{e}_i(\Lambda)}{\widetilde{e}_{i-1}(\Lambda)} \right)^n \le \lambda_{i}^{(n)} \le \binom{d}{{i}} \left(\frac{\widetilde{e}_i(\Lambda)}{\widetilde{e}_{i-1}(\Lambda)} \right)^n
\end{align*}
by using \eqref{eq:Lambda_S}.
Taking the $n$-th root of each value in the above inequality, we get
\begin{align*}
\binom{d}{{i-1}}^{-\frac{1}{n}} \frac{\widetilde{e}_i(\Lambda)}{\widetilde{e}_{i-1}(\Lambda)} \le \big(\lambda_{i}^{(n)}\big)^{\frac{1}{n}} \le \binom{d}{{i}}^{\frac{1}{n}} \frac{\widetilde{e}_i(\Lambda)}{\widetilde{e}_{i-1}(\Lambda)}.
\end{align*}
Hence, we obtain $\big(\lambda_{i}^{(n)}\big)^{\frac{1}{n}} \rightarrow   {\widetilde{e}_i(\Lambda)}/{\widetilde{e}_{i-1}(\Lambda)}$ as $n\rightarrow\infty$.
\end{proof}

\end{document}
